\chapter{用户界面与体验设计}

\section{设计迭代过程}

系统 UI 设计随开发迭代同步演进，经历以下三个阶段。

\textbf{第一阶段——场景驱动草图}：需求分析阶段，团队依据普通用户、家庭管理员、运维人员三类场景，推演各页面的核心信息与操作动线，确定整体布局框架：深色侧边栏承载全局导航，白色卡片区承载数据展示。各功能页面的布局模式在此阶段确定：Dashboard 采用摘要卡片组加图表区的双层结构；校正工作台采用左侧队列加右侧详情的双栏分割；交易列表采用顶部筛选栏加分页表格。

\textbf{第二阶段——设计系统建立}：首个开发迭代中，以 Dashboard 页面为标准件，确定色板、字号层级与间距规范，检验组件复用与色彩层级的一致性，形成贯穿全系统的设计语言（详见下节"设计系统"）。

\textbf{第三阶段——联调反馈与细节调整}：前后端联调过程中，交互细节根据真实操作反馈迭代调整。导入页面的进度展示方式、校正工作台的双栏布局与队列排序逻辑、家庭组织页面的角色权限控件可见性，均在多轮联调中根据反馈完善，最终形成本章所展示的页面形态。

\section{设计系统}

系统采用专业金融工具风格的设计体系，以深色导航与浅色内容区形成清晰的视觉层级。设计系统的核心参数如下：

\begin{table}[H]
\centering
\caption{设计系统色板}
\label{tab:design-system}
\begin{tabular}{p{3.5cm} p{3.5cm} p{3.5cm}}
\toprule
\textbf{色板角色} & \textbf{色值} & \textbf{应用位置} \\
\midrule
深色侧边栏底色 & \texttt{\#06233D} & 左侧导航栏背景 \\
\hline
深色侧边栏悬浮色 & \texttt{\#041A2E} & 导航项悬停/选中态 \\
\hline
品牌主色 & \texttt{\#00A884} & 主要按钮、链接、数据高亮 \\
\hline
页面背景 & \texttt{\#F5F7FA} & 内容区域整体背景 \\
\hline
卡片背景 & \texttt{\#FFFFFF} & 数据卡片、表格、表单容器 \\
\hline
成功色 & \texttt{\#22C55E} & 成功状态、收入标识 \\
\hline
警告色 & \texttt{\#F59E0B} & 警告提示、中等置信度 \\
\hline
错误色 & \texttt{\#EF4444} & 错误状态、支出标识 \\
\hline
信息色 & \texttt{\#3B82F6} & 信息提示、链接、进度条 \\
\bottomrule
\end{tabular}
\end{table}

深色侧边栏提供全局导航，顶部栏显示当前页面标题与用户信息区域。主内容区采用白色卡片承载数据面板、表格与表单，通过规则间距与统一圆角保持视觉一致性。图表采用 ECharts 渲染，配色与设计系统色板保持一致，确保图表与页面风格协调。

\section{主要页面原型}

系统共设计 10 个主要页面，本节逐一描述各页面的功能定位与关键交互。

\subsection{登录与注册页面}

登录与注册页面采用左右分栏布局：左侧为深色品牌展示区，包含应用名称与简短说明；右侧为白色表单卡片，包含用户名与密码输入框，以及登录与注册模式切换按钮。注册模式额外显示邮箱字段。验证成功后，用户进入系统主界面。

\begin{figure}[H]
  \centering
  \uiscreenshot{截图-登陆注册页.png}
  \caption{登录与注册页面截图}
  \label{fig:screenshot-login}
\end{figure}

\subsection{仪表盘页面}

Dashboard 是系统的财务驾驶舱，顶部展示收入、支出、结余、储蓄率四张摘要卡片。中部为支出趋势折线图和分类占比饼图，右侧展示 Top 商户排行榜。用户可通过顶部日期筛选器切换统计周期，从整体趋势快速定位异常支出。

\begin{figure}[H]
  \centering
  \uiscreenshot{截图-dashboard.png}
  \caption{仪表盘页面截图}
  \label{fig:screenshot-dashboard}
\end{figure}

\clearpage
\subsection{账单导入页面}

账单导入页面提供拖拽上传区域，支持微信与支付宝 CSV 账单文件。上传成功后，页面展示最近导入记录列表，包含文件名、平台、状态与处理进度条，支持批量进度轮询。用户可点击导入记录查看详情，或删除历史导入批次。

\begin{figure}[H]
  \centering
  \uiscreenshot{截图-导入账单.png}
  \caption{账单导入页面截图}
  \label{fig:screenshot-import}
\end{figure}

\subsection{交易列表页面}

交易列表页面以筛选栏加分页表格的形式展示所有交易记录。筛选条件包括日期范围、平台（微信/支付宝）、分类、导入文件和关键词搜索。表格每行显示时间、商户、金额、分类与置信度。点击行可展开详情抽屉，展示分类来源、置信度、分类理由与原始交易备注。

\begin{figure}[H]
  \centering
  \uiscreenshot{截图-交易列表.png}
  \caption{交易列表页面截图}
  \label{fig:screenshot-transactions}
\end{figure}

\clearpage
\subsection{分类校正工作台}

校正工作台采用左右分栏布局：左侧为待校正交易队列，按置信度从低到高排序，每项显示商户、金额与当前分类；右侧为当前选中交易的详细信息，包括模型分类建议（来源、置信度、推理理由），以及分类选择控件，允许用户选择新的系统分类或自定义分类并确认修正。确认后交易从队列中移除。

\begin{figure}[H]
  \centering
  \uiscreenshot{截图-分类矫正.png}
  \caption{分类校正工作台截图}
  \label{fig:screenshot-review}
\end{figure}

\subsection{分类管理页面}

分类管理页面展示系统预设分类与用户自定义分类。顶部为分类使用频率统计，按交易数量排序。中部为分类列表，系统分类标记为锁定不可删除，用户自定义分类支持新增、编辑与禁用操作。修改操作通过弹窗表单完成，实时生效。

\begin{figure}[H]
  \centering
  \uiscreenshot{截图-分类管理.png}
  \caption{分类管理页面截图}
  \label{fig:screenshot-categories}
\end{figure}

\clearpage
\subsection{消费人格页面}

消费人格页面由三个区域组成：顶部为 16 型消费人格画像卡片，展示用户所属人格类型及简要描述；中部为四维雷达图（现时偏好、心理账户、炫耀性消费、开放性），展示各维度得分；底部为财务健康评分区域，展示储蓄率、收入稳定性、消费多样性、支出波动性和应急能力的子指标。页面同时提供问卷自评入口，用户可对比主观认知与交易数据画像的偏差。

\begin{figure}[H]
  \centering
  \uiscreenshot{截图-消费人格.png}
  \caption{消费人格页面截图}
  \label{fig:screenshot-personality}
\end{figure}

\subsection{家庭组织页面}

家庭组织页面在用户未加入任何组织时显示创建引导。已有组织时，页面展示三部分内容：左侧为成员列表，每行显示用户名、角色与管理按钮；右上方为家庭财务概览卡片，展示组织合计收入、支出与交易数；右下方为家庭消费人格聚合雷达图，将所有成员交易数据按人格维度聚合展示。家庭所有者和管理员可见角色管理控件与重分类触发按钮。

\begin{figure}[H]
  \centering
  \uiscreenshot{截图-家庭组织.png}
  \caption{家庭组织页面截图}
  \label{fig:screenshot-organization}
\end{figure}

\clearpage
\subsection{报表中心页面}

报表中心页面提供日期范围选择器与导入文件筛选器。用户设置条件后点击生成按钮，系统异步创建报表任务。生成完成后，页面展示报表摘要信息（覆盖时间段、交易数、分类分布），并提供 PDF 下载按钮。历史报表列表按生成时间倒序排列，支持重新下载。

\begin{figure}[H]
  \centering
  \uiscreenshot{截图-报表中心.png}
  \caption{报表中心页面截图}
  \label{fig:screenshot-reports}
\end{figure}

\subsection{系统设置页面}

系统设置页面向所有用户展示当前分类供应商名称、置信度阈值等配置信息。管理员用户额外可见重试队列控制区：包括当前队列状态（等待重试数、重试中数、已失败数），以及“一键重试历史超时”按钮。普通用户仅可见配置展示，管理操作控件隐藏。

\begin{figure}[H]
  \centering
  \uiscreenshot{截图-系统设置.png}
  \caption{系统设置页面截图}
  \label{fig:screenshot-settings}
\end{figure}

\section{界面信息架构与交互组织}

系统界面围绕“全局导航、任务页面、状态反馈、数据可视化”四类信息组织。该划分使用户在不同业务页面之间切换时保持一致的认知模型，也便于后续新增功能时复用相同的交互规则。

\begin{table}[H]
\centering
\caption{界面信息架构与交互组织}
\label{tab:frontend-components}
\begin{tabular}{p{4.5cm} p{9cm}}
\toprule
\textbf{界面区域} & \textbf{设计职责} \\
\midrule
全局导航 & 以深色侧边栏承载主要功能入口，保证用户可在导入、交易、校正、分析、报表等任务之间快速切换 \\
\hline
页面标题与用户区 & 在顶部区域展示当前任务语境、用户身份和必要操作，减少用户对当前位置的判断成本 \\
\hline
任务主区域 & 各页面按“筛选—列表—详情”或“概览—图表—明细”的结构组织内容，匹配财务数据查看和校正流程 \\
\hline
状态反馈 & 导入、报表生成和模型重试等耗时任务通过进度、状态标签和结果摘要展示处理进展 \\
\hline
权限可见性 & 管理员功能、家庭角色管理和普通用户配置展示在界面上区分可见，避免误操作 \\
\hline
数据可视化 & 看板、人格画像和家庭聚合页面采用图表表达趋势、占比和维度评分，辅助用户理解财务状态 \\
\bottomrule
\end{tabular}
\end{table}

交互组织上，系统将身份状态、当前组织和页面筛选条件保持为跨页面可感知的上下文。用户切换页面后不需要重复登录或重新理解导航结构；需要管理员权限的操作仅在具备权限时出现，普通用户界面保持简洁。

\section{用户体验目标}

系统用户体验设计围绕以下四个核心目标展开：

\textbf{减少用户操作步骤}：账单导入采用拖拽上传模式，一次上传即可触发解析→去重→分类的全自动链路。Dashboard 与人格页面的数据通过后端聚合计算，用户无需手动维护分类映射表或公式。批量校正工作台允许用户在单一页面完成所有低置信度交易的复核，无需逐条跳转。

\textbf{分类结果可解释、可校正}：模型分类结果不仅给出类别，还提供置信度与人类可读的推理理由。校正工作台将模型建议与修改控件并置，用户可在了解 AI 推理过程后做出校正决策。校正结果即时生效，无需额外保存步骤。

\textbf{校正数据全链路贯通}：人工校正后的交易分类应实时反映在 Dashboard 聚合、消费人格计算和 PDF 报表中。系统通过统一数据源（transactions 表）确保各分析模块使用一致的最新分类结果，避免出现“已校正但图表未更新”的不一致情况。

\textbf{个人模式不受家庭功能影响}：家庭组织功能采用显式 opt-in 设计——用户需主动创建或加入组织，并在看板和报表页面主动切换至家庭视图。未切换时，所有页面保持个人账本模式，确保现有用户的体验不受影响。
